\scanpage{296}
\mspage{37}
l'égalité n'étant obtenue que si les \unsure{portées|pointures} de $f$
sur les intervalles critiques de $I_{\unsure{\tau|t}}$ \deleted{contenus dans}
qui sont au dessus de $J'$ sont \inserted{toutes} égales à $0$ ou $1$
(donc \inserted{\unsure{ssi}} sur les $\unsure{J|I}\in\mathrm{Cr}(f)^{-1}(\{J\})$, elles sont égales à $1$).

\uline{Cor 4}. Supposons de plus $I'$ \unsure{réduite|réduit}. Alors
\[
\mathrm{Cr}(f)\colon\mathrm{Cr}(\tau,f)\longrightarrow\mathrm{Cr}(\tau')
 \bigl(\simeq\mathrm{Fl}_{\mathrm{tr}}(\tau')\bigr)
\]
est \uline{bijective}.

Donc dans ce cas \inserted{$\tau'$ réduit}, on obtient une application
en sens inverse
\diagram{\begin{tikzcd}[ampersand replacement=\&,column sep=small]
\& \mathrm{Cr}(\tau')\arrow[rr,"\mathrm{Cr}^{\unsure{*|+}}(f)",Rightarrow]\arrow[dr,"\mathrm{Cr}(f)^{-1}"'] \& \& \mathrm{Cr}(\tau)\\
\mathrm{Tr}(\tau')\arrow[ur] \& \& \mathrm{Cr}(\tau,f)\arrow[ur,hook,"\text{inclusion}"'] \& {}
\end{tikzcd}},
application qui transforme \deleted{\unsure{la flèche}}\inserted{l'intervalle}
\deleted{transition}\inserted{critique} $J'$ $(=\{i-1,i\})$ en l'\uline{unique}
intervalle critique $J$ de $I_\tau$ tel que
\[
 f(J)=J'.
\]

% slanted
\margin{\deleted{\unsure{Et pour|pour} \deleted{tout} foncteur $f\colon I_\tau\to I_{\tau'}$, \unsure{pas néc.} croissant ni morphisme d'intervalles,}}
\deleted{Mais on peut aussi \inserted{(sous les \unsure{$=$ hyp. sur $f$})} regarder dans $I_\tau$ l'ens. des
\uline{flèches de transition} \deleted{sur lesquelles}\inserted{ou} (i.e. des
sous-intervalles $J$ de longueur $1$ $\unsure{\{i\mp1,i\}|\{i+1,i\}}$, ou
``\uline{intervalles de transition}''), \inserted{$\unsure{f(u)|f(u_i)}\ne\mathrm{id}$}, \unsure{i.e} telles que
$f\mid J$ non constante. Les \deleted{\unsure{formes}} \unsure{d'où} une
application (supposant $\tau'$ réduit)
\[
\mathrm{Tr}(f)\colon\unsure{\mathrm{Ft}|\mathrm{Tr}}(\tau,f)\longrightarrow\mathrm{Tr}(\tau')
\]
\uline{Cor. 5} : L'image de $\mathrm{Tr}(f)$ est un segment de
l'ens. totalement ordonné $\mathrm{Tr}(\tau')$. Il contient
toutes les flèches de transition \deleted{\unsure{dont le}} de $I_{\tau'}$ dont le}
% archive sheet number: 295
